\documentclass{situationist-preprint}

\title{Institutions for Global Justice}
\msid{POL-13-b}
\msfamily{POL-13}
\msdiscipline{Political Philosophy}
\mstype{Companion manuscript}
\mskeywords{global justice; beneficence; international institutions; assurance; aid}

\begin{document}
\maketitle

\begin{abstract}
POL-13-a argues that the duty of beneficence of ETH-06, to contribute ten per cent of income to the collective task of meeting basic need, applies globally, and that its fulfilment requires institutions that assess need, allocate contributions and assure contributors. This companion asks what such institutions must be like. We argue that the three functions can be performed by different bodies, that assessment and allocation require public standards of accountability but not a world government, and that assurance requires only the publication of aggregate contributions against the compliance floor. We answer the objection that existing aid institutions are too unreliable for the duty to be discharged through them.
\end{abstract}

\section{The three functions}

Definition~2.2 of POL-13-a says that an institution serves the collective task if it assesses need, allocates contributions without duplication or gaps, and assures each contributor that her share belongs to a scheme that reaches the compliance floor. Proposition~3.2 of POL-13-a shows that only institutions can do this globally. It does not say that one institution must.

\begin{proposition}\label{p:division}
The three functions of Definition~2.2 of POL-13-a can be discharged by distinct bodies, provided that the outputs of assessment are public and allocation answers to them.
\end{proposition}

\begin{proof}
Assessment produces information: where basic need lies and what meeting it costs. Allocation uses that information to direct contributions. Assurance reports how far total contributions fall short of, or exceed, what the floor requires. Each function takes the output of the previous one as input and needs nothing else from it. Distinct bodies can therefore perform them if the information passes between them publicly, so that each can be checked against the others.
\end{proof}

The division matters because it lowers what the duty requires of the international order. Assessment is already performed, imperfectly, by statistical agencies and by research on the measurement of poverty \citep{Pogge2002}; allocation is performed by national aid agencies, multilateral funds and private organizations; assurance is the function most conspicuously missing. What is required is not a world government but public standards that bind these bodies to the information produced by assessment, and a public account of contributions against the floor.

\begin{proposition}\label{p:assurance}
Assurance requires only that aggregate contributions be published against the compliance floor of Corollary~3.3 of POL-13-a.
\end{proposition}

\begin{proof}
By that corollary, need is met if contributors holding one tenth of affluent income comply. A contributor is assured when she knows whether this threshold is reached. That is a fact about aggregates, and its publication requires no coercive power over non-contributors.
\end{proof}

Proposition~\ref{p:assurance} distinguishes the present view from one on which the duty is a full-compliance share, which would require institutions able to compel the affluent of every country \citep{Murphy2000}. It also answers \citet{Nagel2005}, who holds that obligations of justice presuppose a shared sovereign: the institutions required here coordinate a humanitarian duty, which Nagel himself accepts holds between all persons, and need not be sovereign.

\begin{objection}
Existing aid institutions are wasteful and sometimes harmful. A duty to contribute through them is a duty to fund their failures.
\end{objection}

\begin{reply}
The duty is to contribute to a task that institutions serve, and Definition~2.2 of POL-13-a states what serving it is. An agency that does not allocate according to public assessment of need does not serve the task, and contributions to it do not discharge the duty. The objection is therefore a reason to direct contributions to bodies that meet the standard, and to build the public assessment that makes the standard applicable, not a reason to think the duty lapses \citep{Singer1972}.
\end{reply}

\section{Conclusion}

The institutions that the global duty of beneficence requires are a division of labour among assessment, allocation and assurance, bound together by public information. They need not be sovereign, and they need secure only the compliance floor. This is what the conclusion of POL-13 requires: a duty that applies globally and is fulfilled through international institutions.

\begin{thebibliography}{9}\small
\bibitem[Murphy(2000)]{Murphy2000} Murphy, L. B. (2000). \emph{Moral Demands in Nonideal Theory}. Oxford University Press.
\bibitem[Nagel(2005)]{Nagel2005} Nagel, T. (2005). The problem of global justice. \emph{Philosophy \& Public Affairs}, 33(2), 113--147.
\bibitem[Pogge(2002)]{Pogge2002} Pogge, T. (2002). \emph{World Poverty and Human Rights}. Polity Press.
\bibitem[Singer(1972)]{Singer1972} Singer, P. (1972). Famine, affluence, and morality. \emph{Philosophy \& Public Affairs}, 1(3), 229--243.
\end{thebibliography}

\end{document}
